\section{Introduction}
\label{sec:intro}

Five labelled aerial scenes of a class the network has never seen. A ResNet-18 trained on $25$ other NWPU-RESISC45 classes labels the query scenes of the new class correctly \num{resisc45:dense_ce:novel_5}\% of the time. Keep $30\%$ of its parameters and the figure falls to \num{resisc45:A_dacis_PMP_30:novel_5}\%. Pick the kept channels at random instead of by any of seven computed scores and it is \num{resisc45:A_random_PMP_30:novel_5}\%. This paper is about that last comparison.

It matters wherever a classifier has to run on hardware that is small and cannot be retrained often: the payload of a small satellite, a drone over a field, a field laptop. On-board inference is already flying~\cite{giuffrida2020cloudscout,giuffrida2022phisat}, the classes that matter on orbit change with the mission, and labels for a new class are few. A pruned network is cheap to run and a prototype head can be updated from five labelled patches without touching the weights, so the combination looks natural. The pruning studies we build on~\cite{li2017pruning,liu2017slimming,molchanov2017pruning,he2017channel} report accuracy on the classes the network was trained for. What survives when the classes at test time are not those the channels were chosen for?

We study that question on ResNet-18. Every method prunes the same architecture: it removes the inner convolution channels of each residual block under one shared depth profile (Sec.~\ref{sec:method}), so two scores differ only in which channels they keep. The scores are random selection, filter $\ell_1$ norm, BatchNorm scale, a Taylor term, and DACIS, which adds three per-channel quantities (gradient norm, activation variance, Fisher ratio over class labels); the three quantities are also run alone. They are trained under seven schedules, among them the PMP schedule of prune, meta-train, prune again that introduced DACIS in an earlier preprint~\cite{uddin2026preprint}, at budgets from $50\%$ to $10\%$ of the parameters. We test on held-out classes of aerial scenes (NWPU-RESISC45~\cite{cheng2017resisc}), of Sentinel-2 patches with all thirteen bands (EuroSAT-MS~\cite{helber2019eurosat}) and of PlantVillage leaf photographs~\cite{hughes2015plantvillage}, with PlantDoc field photographs~\cite{singh2020plantdoc} as a domain-shift check, over \num{resisc45:seeds} random class partitions and with the same test episodes for every method. The first two sets are Earth observation. The crop sets are included because the method under test was first proposed for plant disease; they are ground-level photographs without geospatial context, and we treat them as a secondary setting.

We commit to three claims. First, at $30\%$ of the parameters computed scores do not beat random selection: of \num{cmp:total:PMP} paired comparisons (seven scores, four test settings) \num{cmp:above:PMP} favour the computed score and \num{cmp:below:PMP} favour random selection, and the ranking of the scores changes from one class partition to the next (Sec.~\ref{sec:results}). Second, episodic training on top of cross-entropy pre-training did not raise accuracy on novel classes, so no schedule that contains it beat prune-and-recover. Third, the budget and the length of recovery move accuracy more than the score does, and the saving is smaller in time than in operations: removing \num{cost:macs-red:30}\% of the multiply-accumulates shortened single-thread CPU latency by \num{cost:cpu-red:30}\%. One question stays open. A score computed on the base classes may favour channels that separate those classes and nothing else; a score computed on the target classes would need target labels at deployment, and we did not test it.

The next section places this work among pruning and few-shot studies; Sec.~\ref{sec:method} defines the scores and schedules, Sec.~\ref{sec:setup} the protocol. The notebook that reruns every number in the paper accompanies the submission.
